\paragraph{高係数の同時非負性は非特異領域だけに残る．}
\label{nhd20:reduction}
canonical 表示で $b=-\beta$，$d=-D_0=U+1-S-T$ と置く．
整数 row sums，pair products，正の weights のみを仮定すると，
\[
 D_a,D_b,D_c\ge0
 \quad\Longrightarrow\quad
 b\ge2,\quad d\ge1,\quad \alpha\ge k\ge b+1,
\]
\[
 \Delta=\mathcal C(\mathcal A-\mathcal B)
          +\mathcal B(\mathcal A+1)
 \ge(-\mathcal B)(-\mathcal B-1)>0.
\]
分母の上限，奇数性，clutching，node 条件は不要である．
したがって，階数二を含む全ての特異 Schur 系から
全高 $D$ 非負の模型が生じる可能性はない．
この結論は，特異な正整数局所模型そのものを排除するものではない．

再利用する証明の要点を残す．
\[
 D_a=sd+b-k-1,\quad
 D_b=td-\alpha+k,\quad
 D_c=\alpha-b+1-ud,\qquad
 D_a+D_b+D_c=d
\]
である．まず全 $D$ 零を除く．その場合は
$d=0,\alpha=k=b-1$ で，$b\ge2$ である．
\[
 m=b^2-b+1,\qquad K=s^2+st+t^2-s-t,\qquad
 r=(-s,-t,u)^{\mathsf T},\quad x=(-S,-T,U)^{\mathsf T},
\]
\[
 B_0=\begin{pmatrix}
 1&b-1&-b\\
 -b&1&b-1\\
 b-1&-b&1
 \end{pmatrix}
\]
と置くと，Cramer 恒等式は $\Delta=mK$，$K=m-1$，
$x=r-B_0r/m$ を与える．整数性から各成分で $m\mid B_0r$ である一方，
\[
 B_0^{\mathsf T}B_0=m(3I-J),\qquad
 \lVert r\rVert^2=2m-1,\qquad \mathbf1^{\mathsf T}r=-1
\]
である．ここで $J$ は全成分が一の三次行列である．従って
$m^2\mid\lVert B_0r\rVert^2=m(6m-4)$，従って $m\mid4$ となる．
$m$ は三以上の奇整数だから矛盾する．
特に $D_0=0$ なら高係数には正と負がともに現れる．

同時非負性には $d\ge1$ が必要で，
$0\le D_a\le d$ から
\[
 (s-1)d+b-1\le k\le sd+b-1
\]
を得る．$k=b-1$ なら $s=1$，$u=t$ となり，
線形系の第三式は
\[
 t(b-1)T-tbU=b(b-1)>0
\]
となる．しかし $U=S+T+d-1>T$ なので左辺は負である．
$k=b$ なら $(s-1)d\le1$ であり，$s=1$ または $(s,d)=(2,1)$ となる．
前者では第三式 $tb(T-U)=b^2-1>0$ が同じ理由で矛盾する．
後者では $\alpha=b+t$，$U=S+T$ であり，第一式と第三式の線形結合が
\[
 t(1+b^2+tb)S=t+2b-b^3<t
\]
を要求するため，$S\ge1$ に反する．
この境界排除は determinant で割らないので，途中で特異系を落としていない．

以上から $k\ge b+1$ であり，
$D_c\ge0$ と $k\le sd+b-1$ から
$\alpha-k\ge(t-1)d\ge0$ が従う．
$\bar b=-\mathcal B=s+tb$ と置けば
\[
 \mathcal C-\bar b=(s-1)b+u(k-b)\ge0,\qquad
 \Delta=\bar b(\bar b-1)+(\mathcal C-\bar b)(\mathcal A+\bar b)
\]
により determinant の狭義正値性を得る．

残る非特異領域での一般符号問題は未解決である．
この領域の整数局所模型自体は存在する．例えば
\[
 (s,t,\alpha,\beta,k)=(1,2,10,-2,5),\qquad
 (S,T,U)=(4,5,11)
\]
は $q=(5,11,23,1)$，$D=(-1,1,3,-3)$，
正の weights $(5,109/6,45/22,1)$ を与える．
従って $\alpha\ge k\ge b+1$ と $\Delta>0$ だけで局所模型を捨てず，
Cramer 整除性と全高 $D$ の非負性を引き続き課す必要がある．
\nolinkurl{verify_high_d_boundary_v20.py} は
四十二の多項式恒等式と三つの厳密行列回帰を検算する．
一般命題の根拠は上の恒等式・不等式であり，有限範囲の代入結果ではない．
